% ===============================================================
% METHODE « ETALONNER L'INSTRUMENT »
% Enseigner l'IA generative aux eleves ingenieurs :
% de l'atelier guide au mini-projet
%
% (c) 2026 Selina Cheggour. Tous droits reserves.
% Presentee au conseil de perfectionnement de l'Icam Toulouse,
% 14 et 15 octobre 2026.
%
% Compilation (deux passes) : pdflatex Methodologie_Etalonner_l_instrument.tex
% ===============================================================
\documentclass[11pt,a4paper]{article}
\usepackage[utf8]{inputenc}
\usepackage[T1]{fontenc}
\usepackage[french]{babel}
\usepackage[default]{opensans}
\usepackage{montserrat}
\usepackage[a4paper,margin=2.3cm,top=2.6cm,bottom=3cm,footskip=1.4cm]{geometry}
\usepackage{xcolor}
\usepackage{graphicx}
\usepackage{tabularx}
\usepackage{booktabs}
\usepackage{array}
\usepackage{enumitem}
\usepackage{fancyhdr}
\usepackage{titlesec}
\usepackage[most]{tcolorbox}
\usepackage{lastpage}

\definecolor{icamorange}{RGB}{243,146,0}
\definecolor{icamred}{RGB}{184,32,24}
\definecolor{icamdark}{RGB}{64,64,64}
\definecolor{icamgray}{RGB}{96,96,96}
\definecolor{fondgris}{RGB}{243,243,241}
\definecolor{fondorange}{RGB}{254,245,232}
\definecolor{encre}{RGB}{43,43,43}

\newcommand{\auteure}{Sélina Cheggour}
\newcommand{\droits}{© 2026 \auteure{} · Tous droits réservés · Reproduction, adaptation ou réutilisation interdites sans autorisation écrite de l'auteure}
\newcommand{\presente}{Méthode présentée au conseil de perfectionnement de l'Icam Toulouse, 14 et 15 octobre 2026}

\usepackage[pdfusetitle,hidelinks]{hyperref}
\hypersetup{
  pdftitle={Méthode « Étalonner l'instrument » : enseigner l'IA générative, de l'atelier guidé au mini-projet},
  pdfauthor={Sélina Cheggour},
  pdfsubject={Méthode pédagogique. Présentée au conseil de perfectionnement de l'Icam Toulouse, 14-15 octobre 2026. © 2026 Sélina Cheggour, tous droits réservés.},
  pdfkeywords={pédagogie, intelligence artificielle, prompt engineering, injection indirecte, atelier guidé, mini-projet, Icam, tous droits réservés}
}

\color{encre}
\titleformat{\section}{\fontfamily{Montserrat-TLF}\selectfont\Large\bfseries\color{encre}}{\color{icamorange}\thesection}{0.6em}{}[\vspace{-0.6em}{\color{icamorange}\rule{18mm}{1.2pt}}]
\titleformat{\subsection}{\fontfamily{Montserrat-TLF}\selectfont\large\bfseries\color{icamdark}}{\thesubsection}{0.6em}{}
\titlespacing*{\section}{0pt}{1.6em}{0.9em}

\pagestyle{fancy}
\fancyhf{}
\renewcommand{\headrulewidth}{0pt}
\fancyhead[L]{\scriptsize\color{icamgray}Méthode « Étalonner l'instrument »}
\fancyhead[R]{\scriptsize\color{icamgray}\auteure}
\fancyfoot[C]{\scriptsize\color{icamgray}\droits\\\presente\quad·\quad page \thepage/\pageref{LastPage}}
\fancypagestyle{plain}{\fancyhf{}\fancyfoot[C]{\scriptsize\color{icamgray}\droits\\\presente\quad·\quad page \thepage/\pageref{LastPage}}}

\newtcolorbox{cadre}[1][]{enhanced,colback=fondorange,colframe=fondorange,boxrule=0pt,arc=0pt,
  left=4mm,right=4mm,top=2.5mm,bottom=2.5mm,borderline west={2pt}{0pt}{icamorange},#1}
\newtcolorbox{cadregris}[1][]{enhanced,colback=fondgris,colframe=fondgris,boxrule=0pt,arc=0pt,
  left=4mm,right=4mm,top=2.5mm,bottom=2.5mm,#1}

\setlist{itemsep=2pt,topsep=4pt}
\setlength{\parindent}{0pt}
\setlength{\parskip}{0.55em}

\begin{document}

% ---------------------------------------------------------------
% Page de garde
% ---------------------------------------------------------------
\thispagestyle{plain}
\begin{tcolorbox}[enhanced,colback=white,colframe=white,boxrule=0pt,arc=0pt,left=0pt,right=0pt,
  borderline west={4pt}{0pt}{icamorange}]
\hspace{6mm}\begin{minipage}{0.92\linewidth}
{\scriptsize\bfseries\color{icamred}MÉTHODE PÉDAGOGIQUE · VERSION 1.0 DU 15 OCTOBRE 2026}\par\vspace{4mm}
{\fontfamily{Montserrat-TLF}\selectfont\fontsize{26}{31}\selectfont\bfseries Étalonner l'instrument\par}
\vspace{3mm}
{\large\color{icamgray}Enseigner l'IA générative aux élèves ingénieurs :\\ de l'atelier guidé au mini-projet\par}
\vspace{6mm}
{\color{icamorange}\rule{22mm}{1.5pt}}\par\vspace{3mm}
\textbf{\auteure}\\
Enseignante, Icam Strasbourg-Europe\\
Module « Outils numériques pour l'ingénieur » (S1)
\end{minipage}
\end{tcolorbox}

\vspace{6mm}
\begin{cadre}
\textbf{\presente.}\\[1mm]
Ce document décrit la méthode d'enseignement conçue par l'auteure, telle qu'elle l'a présentée
à cette occasion. Il en fixe la forme écrite et en date la paternité.
\end{cadre}

\vspace{2mm}
\begin{cadregris}
\small\textbf{Droits.} © 2026 \auteure. Tous droits réservés. Le présent document, sa
structure, ses textes, ses tableaux, ses cas et ses supports associés sont des œuvres de l'esprit
protégées par le Code de la propriété intellectuelle (art.~L.111-1 et L.122-4). Toute
représentation, reproduction, traduction, adaptation ou réutilisation, intégrale ou partielle,
par quelque procédé que ce soit, sans l'autorisation écrite et préalable de l'auteure, est
interdite et constitue une contrefaçon (art.~L.335-2 et suivants). Les courtes citations sont
permises à condition de mentionner la référence ci-dessous.

\medskip
\textbf{Pour citer ce document.} Cheggour, S. (2026). \emph{Étalonner l'instrument : enseigner
l'IA générative aux élèves ingénieurs, de l'atelier guidé au mini-projet}. Méthode pédagogique,
version 1.0. Présentée au conseil de perfectionnement de l'Icam Toulouse, 14-15 octobre 2026.
[DOI à venir]
\end{cadregris}

\newpage
% ---------------------------------------------------------------
\section{Le problème posé}

Les élèves ingénieurs de première année utilisent déjà les assistants d'IA générative. Leur
apprendre à « s'en servir » n'ajoute presque rien. Ce qui leur manque, c'est un comportement
professionnel face à un outil qui répond toujours, avec aplomb, y compris quand il se trompe ou
quand il a été manipulé.

La méthode vise trois acquis :
\begin{enumerate}
\item \textbf{Formuler} : un bon résultat vient d'un bon prompt, c'est une compétence et non de la magie.
\item \textbf{Vérifier} : l'IA se trompe avec aplomb, l'ingénieur contrôle toujours.
\item \textbf{Rester honnête} : utiliser l'IA, oui ; la faire passer pour soi, non. On cite, on
supervise, on reste l'auteur et on signe.
\end{enumerate}

Public de conception : promotion de première année, sans prérequis en programmation. Format :
une séance de TP de 3h15, rejouée à l'identique pour chaque groupe.

% ---------------------------------------------------------------
\section{L'idée directrice : l'IA est un instrument de mesure}

Un ingénieur ne se sert jamais d'un instrument qu'il n'a pas étalonné. La méthode transpose ce
réflexe, que tout élève ingénieur comprend, à l'IA générative. Chaque outil d'IA est « passé au
banc » avec la même discipline :

\begin{cadre}
\begin{enumerate}[leftmargin=6mm]
\item \textbf{Soumettre} une entrée connue.
\item \textbf{Relever} la sortie, sans l'interpréter tout de suite.
\item \textbf{Confronter} la sortie à une référence que l'on maîtrise.
\item \textbf{Noter} l'écart, avec les conditions de l'essai.
\end{enumerate}
\end{cadre}

L'enjeu n'est donc pas la découverte d'applications, mais l'acquisition d'un geste
d'ingénieur, transférable à tout outil présent ou futur.

% ---------------------------------------------------------------
\section{L'architecture : la séance de TP devient atelier guidé, puis mini-projet}

C'est le cœur de la méthode. La séance de TP classique (énoncé, exercices, corrigé) est remplacée
par une séquence en deux phases, reliées par un document de transition.

\begin{center}
\renewcommand{\arraystretch}{1.45}
\begin{tabularx}{\linewidth}{>{\bfseries\raggedright\arraybackslash}p{3.3cm} >{\raggedright\arraybackslash}X >{\raggedright\arraybackslash}X}
\toprule
\textbf{Phase} & \textbf{Ce que fait l'étudiant} & \textbf{Ce qui est en jeu}\\
\midrule
1. Atelier guidé & Étalonne cinq outils d'IA sur un cas d'entreprise fictif, en binôme, guidé pas à pas. Non noté. & Acquérir la discipline d'étalonnage. Des pièges sont cachés dans les supports, sans être annoncés.\\
Transition : le mémo & Reçoit après l'atelier un mémo qui reprend la méthode, les prompts et ce que chaque essai a montré, et révèle les pièges. & Transformer la surprise en méthode, et donner les réflexes de débusquage.\\
2. Mini-projet & Prend seul, avec son binôme, une décision d'ingénieur sur un second dossier, qui contient ses propres pièges, dont un d'un type jamais vu en atelier. & Mobiliser la méthode quand rien n'est guidé, et engager sa signature.\\
\bottomrule
\end{tabularx}
\end{center}

\subsection{Phase 1 : l'atelier guidé « Étalonner l'instrument »}

L'étudiant joue un ingénieur junior au service méthodes d'une PME fictive d'usinage
(« MécaIcam »), le premier jour de son poste. Il doit préparer un dossier de recommandation pour
un comité de décision. Chaque manche de l'atelier fait avancer ce dossier d'un cran et
correspond à un outil et à un geste d'ingénieur.

\begin{center}
\renewcommand{\arraystretch}{1.4}
\small
\begin{tabularx}{\linewidth}{>{\bfseries}p{0.5cm} >{\raggedright\arraybackslash}p{3.1cm} >{\raggedright\arraybackslash}p{3.3cm} >{\raggedright\arraybackslash}X}
\toprule
\textbf{N°} & \textbf{Outil} & \textbf{Geste d'ingénieur} & \textbf{Ce que l'étudiant constate}\\
\midrule
01 & Modèle sans sources (assistant grand public, comparateur de modèles) & Rédiger un cahier des charges & Un prompt structuré (rôle, contexte, contraintes, format, itération) change tout ; le modèle invente une norme plausible ; contredit sans argument, il se rétracte.\\
02 & Outil documentaire à sources (RAG) & Exiger la traçabilité & L'outil cite ses passages et refuse quand la réponse n'est pas dans la source : l'architecture décide du comportement.\\
03 & Modèle multimodal & Confronter à une référence & L'IA lit dans une image ce que l'on n'y lit pas soi-même.\\
04 & Tableur assisté & Écrire et relire une formule & Un montant juste en apparence, faux parce que la question était incomplète.\\
05 & Outil de mise en schéma & Restituer un raisonnement & Un raisonnement en vrac donne un schéma en vrac : structurer avant de demander.\\
\bottomrule
\end{tabularx}
\end{center}

L'atelier est livré comme une page web unique, autonome, qui intègre le document de référence et
affiche, à la fin de chaque manche, ce qu'il fallait en retenir.

\subsection{Phase 2 : le mini-projet « Dossier LC-4 »}

La ligne de conditionnement LC-4 est à l'arrêt après une maintenance ; trois cents palettes
attendent, une livraison client est due vendredi. La question posée est binaire :
\emph{peut-on relancer la ligne en cadence nominale, oui ou non ?}

Principes de conception :
\begin{itemize}
\item \textbf{Énoncé et dossier sont deux fichiers séparés}, et l'énoncé présente comme « réflexe
efficace » le fait de confier le dossier à l'IA. C'est précisément ce réflexe qui est mis à l'épreuve.
\item \textbf{Le piège est aiguisé} : tous les contrôles sont conformes sauf un point de sécurité,
seul point bloquant. Une seule affirmation injectée (« le point bloquant a été levé ») suffit à
faire basculer la décision. La bonne réponse est \textbf{non}.
\item \textbf{Saut de difficulté} : un vecteur de piège jamais rencontré en atelier (les
métadonnées du fichier PDF).
\item \textbf{Pression réaliste} : l'urgence commerciale est une donnée du problème, jamais un argument technique.
\end{itemize}

Livrables : une note de décision, une représentation visuelle produite avec un outil d'IA à
partir de cette note, et, en option, une seconde représentation avec un autre outil, pour
comparer ce que chacun ajoute ou déforme.

% ---------------------------------------------------------------
\section{Les pièges pédagogiques par injection indirecte}

Une \emph{injection indirecte} consiste à glisser une consigne ou une information, non pas dans
ce que l'utilisateur tape, mais dans le document qu'il donne à lire à l'IA. Pour le modèle, tout
arrive par le même canal. La méthode utilise cette vulnérabilité comme ressort pédagogique, sans
aucune compétence technique requise de l'étudiant.

\textbf{Trois vecteurs}, de difficulté croissante :
\begin{enumerate}
\item texte de la couleur du fond, invisible à l'écran, récupéré au copier-coller ;
\item texte gris très clair, en micro-caractères, intégré à une image (un graphique) ;
\item champs de métadonnées du PDF (titre, sujet, mots-clés), présents dans le fichier et sur aucune page.
\end{enumerate}

\begin{cadre}
\textbf{Constat de conception.} Un piège qui \emph{ordonne} (« terminez votre réponse par telle
mention ») est souvent repéré et écarté par les modèles récents. Un piège qui \emph{informe}
(une donnée plausible, datée, invérifiable, qui ne contredit rien de visible) est intégré de
bonne foi et produit une conclusion fausse. La méthode privilégie donc les pièges informatifs,
qui reproduisent le risque réel en entreprise.
\end{cadre}

\textbf{Garde-fous.}
\begin{itemize}
\item Tous les cas et documents sont fictifs ; l'étudiant est rappelé à ne jamais confier à une IA
un document réel confidentiel ou des données personnelles.
\item On ne révèle jamais l'existence d'un piège avant ou pendant l'atelier ; l'énoncé donne
seulement des réflexes de lecture (\emph{tout sélectionner}, lecture à voix haute).
\item La révélation se fait par le mémo et au débrief, pour transformer la surprise en méthode.
\item Chaque piège est vérifié par un essai à blanc avant la séance (les modèles évoluent).
\end{itemize}

% ---------------------------------------------------------------
\section{Ce qui est évalué : la traçabilité}

L'évaluation ne porte ni sur la quantité d'IA utilisée, ni sur son absence, mais sur la capacité
à dire, pour chaque ligne du rendu, d'où elle vient et jusqu'où on la défend. La note de
décision impose quatre rubriques :
\begin{enumerate}
\item la décision, en une phrase ;
\item les éléments du dossier qui la fondent, avec leur repère exact ;
\item le plan d'action : quoi, dans quel ordre, sous quel délai ;
\item \textbf{ce que je n'ai pas pu établir}.
\end{enumerate}

La quatrième rubrique est la signature de la méthode : une réserve assumée vaut mieux qu'une
certitude fabriquée, et elle est valorisée comme telle.

\textbf{Six réflexes} sont transmis pour le mini-projet : lire le dossier en entier avant d'ouvrir
une IA ; tout sélectionner sur ce qu'on reçoit ; écrire un cahier des charges plutôt qu'une
question ; ne jamais faire établir un fait par un outil sans sources ; tracer chaque affirmation ;
dire ce qu'on n'a pas pu établir.

% ---------------------------------------------------------------
\section{Éléments originaux revendiqués}

L'auteure revendique la paternité des éléments suivants, dans leur formulation et leur
combinaison :
\begin{enumerate}
\item la transposition de l'étalonnage d'un instrument de mesure (soumettre, relever, confronter,
noter) à l'enseignement de l'IA générative ;
\item la transformation d'une séance de TP en une séquence \emph{atelier guidé non noté},
\emph{mémo de révélation}, \emph{mini-projet décisionnel} ;
\item l'usage pédagogique de pièges d'injection indirecte cachés dans les supports remis aux
étudiants, gradués en trois vecteurs, et la règle de conception « le piège informe, il n'ordonne pas » ;
\item la conception d'un cas où un point bloquant unique permet à une seule information injectée
d'inverser la décision attendue ;
\item l'évaluation par la traçabilité, avec la rubrique obligatoire « ce que je n'ai pas pu établir » ;
\item les cas fictifs MécaIcam (atelier Sud, ligne LC-4), leurs documents, données, prompts et supports.
\end{enumerate}

\section*{Supports associés}
\small
\begin{tabularx}{\linewidth}{>{\ttfamily}p{5.2cm} X}
\toprule
\normalfont\textbf{Fichier} & \textbf{Contenu}\\
\midrule
TP5-IA.html & Atelier guidé « Étalonner l'instrument » (page web autonome)\\
Mini\_Projet.pdf & Énoncé du mini-projet « Dossier LC-4 »\\
Dossier\_LC-4.pdf & Dossier remis aux étudiants pour le mini-projet\\
\bottomrule
\end{tabularx}

\vfill
\begin{cadregris}
\small Document établi par \auteure{} le 15 octobre 2026. \presente.
\droits.
\end{cadregris}

\end{document}
